\nuevaReceta{Leche con Canela y Limón}{aprox. 500 ml · 2-3 vasos}{10 min activos · 2 h 30 min total}{receta:leche_canela_limon}

    \begin{minipage}[t]{0.35\textwidth}
        \section*{Ingredientes}
        \begin{itemize}[leftmargin=*]
            \item 1 l de leche entera
            \item 1 rama de canela
            \item la piel de 1 limón
            \item 60 g de azúcar, aprox.
            \item canela molida (opcional)
        \end{itemize}
    \end{minipage}
    \hfill
    \begin{minipage}[t]{0.6\textwidth}
        \section*{Preparación}
        \begin{enumerate}
            \item \textbf{Preparar el limón:} Lavarlo bien y sacar la piel en tiras, procurando no llevarse la parte blanca para que no amargue.
            \item \textbf{Calentar:} Poner en un cazo la leche, la rama de canela, la piel de limón y el azúcar. Calentar a fuego medio, removiendo hasta que el azúcar se disuelva.
            \item \textbf{Infusionar:} Cuando la leche esté a punto de hervir, retirar del fuego, tapar el cazo y dejar reposar durante 20-30 minutos.
            \item \textbf{Colar:} Retirar la canela y la piel de limón. Probar y ajustar el azúcar si fuera necesario.
            \item \textbf{Enfriar:} Dejar que pierda temperatura y guardar en la nevera durante al menos 2 horas.
            \item \textbf{Servir:} Remover antes de repartirla en vasos y, si se desea, terminar con una pizca de canela molida.
        \end{enumerate}
    \end{minipage}

    \vspace{1em}
    \begin{tcolorbox}[colback=orange!5, colframe=orange!50, title=Consejo, size=small]
        También se puede tomar templada. Para un sabor más intenso, deja la canela y el limón dentro mientras se enfría y cuélala justo antes de guardarla en la nevera.

        \medskip
        \textbf{Rendimiento real:} Según la experiencia de casa, por cada litro de leche utilizado suelen quedar aproximadamente 500 ml, ya que durante la preparación se evapora una cantidad considerable.
    \end{tcolorbox}
\end{receta}
